\documentclass[11pt]{article}
\usepackage[margin=0.7in]{geometry}
\usepackage{amsmath,amssymb}
\usepackage{enumitem}
\usepackage{tikz}
\usepackage{pgfplots}
\pgfplotsset{compat=1.18}
\usepackage{fancyhdr}
\usepackage{parskip}
\setlist[enumerate]{leftmargin=*, itemsep=6pt}
\pagestyle{fancy}
\fancyhf{}
\fancyfoot[L]{\small\textit{en}Vision\textsuperscript{\textregistered} Algebra 2}
\fancyfoot[C]{\small 6-3 Lesson Quiz, short form}
\fancyfoot[R]{\small\thepage}
\renewcommand{\headrulewidth}{0pt}
\newcommand{\choice}[1]{\tikz[baseline=-0.6ex]\node[draw, circle, inner sep=0.6pt, minimum size=3.3mm, font=\scriptsize\bfseries] {#1};\;}
\newcommand{\checkbox}{$\square$\;}

\begin{document}

\noindent\textbf{Name}\; \rule{6cm}{0.4pt} \hfill Period \rule{1.2cm}{0.4pt} \quad Date \rule{2.2cm}{0.4pt}

\vspace{4pt}\noindent\rule{\textwidth}{0.6pt}

\vspace{2pt}\noindent{\LARGE\textbf{6-3 \; Lesson Quiz: Logarithms}} \hfill {\small short form}

\vspace{4pt}\noindent\rule{\textwidth}{0.6pt}

\noindent{\small Questions keep their Savvas Lesson Quiz numbers (1, 3, 5), the numbers your packet uses.}

\vspace{6pt}
\textbf{1.} Which function is the inverse of the exponential function $y=3^x$?

\vspace{4pt}
\begin{itemize}[leftmargin=1.2em, itemsep=5pt]
\item[] \choice{A} $y=x^3$ \qquad\qquad \choice{B} $y=3^x$ \qquad\qquad \choice{C} $y=\log_3 x$ \qquad\qquad \choice{D} $y=\log_x 3$
\end{itemize}

\vspace{6pt}
Explain your choice. Use the words \emph{base}, \emph{exponent} and \emph{inverse}.

\vspace{12pt}\rule{\linewidth}{0.4pt}

\vspace{12pt}\rule{\linewidth}{0.4pt}

\vspace{10pt}\noindent\rule{\textwidth}{0.4pt}

\vspace{6pt}
\textbf{3.} Which is the logarithmic equation $\log_8 64=a$, written in exponential form?

\vspace{4pt}
\begin{itemize}[leftmargin=1.2em, itemsep=5pt]
\item[] \choice{A} $64^a=8$ \qquad\qquad \choice{B} $a^8=64$ \qquad\qquad \choice{C} $a^{64}=8$ \qquad\qquad \choice{D} $8^a=64$
\end{itemize}

\vspace{10pt}\noindent\rule{\textwidth}{0.4pt}

\vspace{6pt}
\textbf{5.} What is the solution to the equation $\log(2x+4)=2$? Round to the nearest thousandth, if necessary.

\vspace{4pt}
Show your work.

\vspace{14pt}\rule{0.6\linewidth}{0.4pt}

\vspace{14pt}\rule{0.6\linewidth}{0.4pt}

\vspace{14pt}\rule{0.6\linewidth}{0.4pt}

\vspace{14pt}\rule{0.6\linewidth}{0.4pt}

\vspace{12pt}
$x=$ \rule{2.5cm}{0.4pt}

\vfill
\noindent\textbf{Scoring:} Question 1 and Question 5 are 10 points each. Question 3 is a bonus of 5 points. Total out of 20.

\newpage
\noindent{\LARGE\textbf{6-3 Lesson Quiz, short form: answer key}}

\vspace{4pt}\noindent\rule{\textwidth}{0.6pt}

\begin{enumerate}
\item[\textbf{1.}] \textbf{C, $y=\log_3 x$.} The logarithmic function $y=\log_b x$ is the inverse of the exponential function $y=b^x$; here the base is 3. ($\log_3 x$ is the exponent you raise the base 3 to in order to get $x$, so it undoes $y=3^x$.) A is a power function (the variable is the base, not the exponent); B is the function itself; D swaps the base and the argument. Savvas quiz item 1, DOK 2, skill F69. Example 1 (continued), Try It 1 and the packet Practice. \textbf{10 points:} 7 for C; 3 for an explanation that names the base 3 and says the logarithm undoes (is the inverse of) raising 3 to an exponent.
\item[\textbf{3.}] \textbf{D, $8^a=64$.} The base 8 raised to the exponent $a$ gives the argument 64 (and $a=2$). Distractor A swaps base and argument; B and C put $a$ in the base. Savvas quiz item 3, DOK 2, skill F73. Taught by the form switch in Example 5 Part B and Try It 5a (and the packet's top box). \textbf{Bonus 5:} D.
\item[\textbf{5.}] \textbf{$x=48$.} Exponential form: $10^2=2x+4$, so $100=2x+4$, $2x=96$, $x=48$. Check: $\log(2\cdot48+4)=\log100=2$. Savvas quiz item 5, DOK 2, skill F69. Example 5 Part B and Try It 5a. \textbf{10 points:} 4 for the exponential form $10^2=2x+4$ (or $100=2x+4$); 6 for $x=48$. A bare 48 with no exponential form earns 6 of 10.
\end{enumerate}

\textbf{Items left off this form} (available for a retake or Form B): Savvas item 2, the best estimate of $\ln16$ (answer B, 2.77; DOK 2, F73; Example 4, cut; $\ln16\approx2.7726$ by calculator); Savvas item 4, select all the logarithms evaluated correctly (answer A and E: $\ln1=0$, $\log_5\frac1{125}=-3$; DOK 1, F69; Example 3, cut).

\textbf{IXL replacement:} SmartScore 70 or above on 8RK and BXU replaces this quiz with 20/20.

\end{document}
